\chapter{Centred Reciprocal Geometry, U(1) Holonomy, and Hubble Normalization}
\label{ch:inverse-holonomy}

\section{Scope and epistemic boundary}
This chapter adds a new layer to the critical-axis programme while preserving the v0.1 claim boundary.  Three ingredients are considered together:
\begin{enumerate}
\item the centred reciprocal chart generated by $u=s-\tfrac12$ and $u\mapsto1/u$;
\item $U(1)$ loop holonomy, with the Aharonov--Bohm phase as the reference gauge-theory realization;
\item a Hubble inverse-time scale $H_0$, used only after nondimensionalization as a candidate global radial normalization.
\end{enumerate}
The broader TIR object $W_{ij}$ may be used elsewhere as a holonomic relation/information connector, but its definition and dynamics are explicitly out of scope here and are deferred to the parent Theory of Informational Relations.

No statement in this chapter identifies the electromagnetic Aharonov--Bohm potential with cosmological expansion as an established physical equality.  Their use as two sectors of one ``potential-tension'' architecture is a \emph{model postulate}.  Likewise, the Hubble radius $c/H_0$ is used as a normalization scale; in a general cosmology it is not automatically a causal or event horizon.

\section{Centre the critical strip before inversion}
Write
\begin{equation}
u=s-\frac12.
\label{eq:centered-u}
\end{equation}
Then the zeta anti-linear involution
\begin{equation}
J(s)=1-\overline{s}
\end{equation}
becomes
\begin{equation}
\widetilde J(u)=-\overline u.
\label{eq:centered-J}
\end{equation}

\begin{proposition}[Centred fixed locus]
\label{prop:centered-fixed-locus}
The fixed set of $\widetilde J$ is the imaginary axis $\Ree u=0$, which is exactly the critical line $\Ree s=1/2$ in the original coordinate.
\end{proposition}

\begin{proof}
The condition $-\overline u=u$ is equivalent to $\Ree u=0$.  Equation \eqref{eq:centered-u} then gives $\Ree s=1/2$.
\end{proof}

This translation is elementary, but it removes an ambiguity that is important for the geometric discussion: the balanced value $1/2$ is represented by the origin of the transverse coordinate.

\section{The reciprocal chart}
On the punctured centred plane define
\begin{equation}
\mathcal I(u)=\frac1u,
\qquad u\neq0.
\label{eq:reciprocal}
\end{equation}
The reciprocal map is holomorphic and conformal on $\mathbb C^\times$.  It is therefore important not to call $1/u$ itself a non-conformal compactification.  The non-conformal step needed to place infinity on a finite boundary is introduced separately below.

\begin{proposition}[Reciprocal involution]
\label{prop:reciprocal-involution}
For every $u\neq0$,
\begin{equation}
\mathcal I(\mathcal I(u))=u.
\end{equation}
\end{proposition}

\begin{proof}
Immediate from $1/(1/u)=u$.
\end{proof}

\begin{proposition}[Compatibility with zeta complement]
\label{prop:reciprocal-commutes}
The centred reciprocal map commutes with the centred zeta involution:
\begin{equation}
\mathcal I\circ\widetilde J
=
\widetilde J\circ\mathcal I
\qquad\text{on }\mathbb C^\times.
\label{eq:commutation}
\end{equation}
\end{proposition}

\begin{proof}
For $u\neq0$,
\begin{equation}
\mathcal I(\widetilde J(u))
=
\frac1{-\overline u}
=
-\frac1{\overline u}
=
-\overline{\frac1u}
=
\widetilde J(\mathcal I(u)).
\end{equation}
\end{proof}

\begin{corollary}[Critical-axis invariance]
\label{cor:critical-axis-inversion}
The critical axis is invariant as a set under reciprocal inversion.  In centred coordinates,
\begin{equation}
u=i\tau
\quad\Longrightarrow\quad
\mathcal I(u)
=
-\frac{i}{\tau},
\qquad \tau\neq0.
\label{eq:critical-inverse}
\end{equation}
Thus the reciprocal chart performs the exact duality
\begin{equation}
\tau\longleftrightarrow-\frac1\tau
\end{equation}
within the critical axis.
\end{corollary}

This is an exact theorem independent of the Riemann hypothesis.  It does \emph{not} say that a zeta zero is mapped to another zeta zero; that would require an additional functional identity not established here.

\section{Finite disk representation: centre and boundary}
To represent the centred point as the centre of a finite disk and the asymptotic region as its boundary, define
\begin{equation}
\mathcal C(u)=\frac{u}{1+|u|}.
\label{eq:radial-compactification}
\end{equation}
This map is intentionally non-holomorphic because of $|u|$.  Its radius is
\begin{equation}
\rho(u)=\frac{|u|}{1+|u|}.
\label{eq:compact-radius}
\end{equation}

\begin{proposition}[Centre--boundary compactification]
\label{prop:center-boundary}
The map $\mathcal C$ satisfies
\begin{equation}
\mathcal C(0)=0,
\qquad
|\mathcal C(u)|<1,
\qquad
|u|\to\infty\Longrightarrow|\mathcal C(u)|\to1.
\end{equation}
\end{proposition}

Hence the geometry requested for the model is implemented by two distinct operations:
\begin{equation}
s
\longmapsto
u=s-\frac12
\longmapsto
\left\{
\begin{array}{ll}
\mathcal I(u)=1/u, & \text{reciprocal/conformal chart},\\[4pt]
\mathcal C(u)=u/(1+|u|), & \text{finite non-conformal disk chart}.
\end{array}
\right.
\end{equation}
Keeping these operations separate prevents a category error: reciprocal inversion and radial compactification do different mathematical jobs.

\section{What can be promoted about $\sigma=\Ree s$}
The v0.1 ledger treated the statement ``$\sigma=\Ree s$ is the correct probability coordinate'' as a single model postulate.  Part of that statement can be sharpened.

Let $x=\Ree s\in[0,1]$ be the transverse critical-strip coordinate and let $p(x)$ be an affine map to a binary probability coordinate.  Impose only the endpoint identifications
\begin{equation}
p(0)=0,
\qquad
p(1)=1.
\label{eq:endpoints}
\end{equation}

\begin{theorem}[Unique affine strip--simplex coordinate]
\label{thm:affine-strip-simplex}
If $p(x)=ax+b$ is affine and satisfies \eqref{eq:endpoints}, then
\begin{equation}
p(x)=x.
\end{equation}
Consequently
\begin{equation}
p(1-x)=1-p(x),
\end{equation}
so the zeta complement $x\mapsto1-x$ is represented exactly by binary probability complement.
\end{theorem}

\begin{proof}
From $p(0)=0$ one obtains $b=0$, and from $p(1)=1$ one obtains $a=1$.
\end{proof}

This promotes the \emph{coordinate-theoretic} part of the old postulate: under the explicit affine and endpoint-preserving assumptions, the choice $p=\Ree s$ is unique.  What remains open is the physically and analytically stronger statement that the zero-state representation of $\xi$ is compelled to use that affine binary coordinate as an amplitude population.  That remainder stays a model postulate/open representation debt.

\section{Aharonov--Bohm phase as a reference $U(1)$ holonomy}
For a $U(1)$ connection, a closed loop $\gamma$ carries the gauge-invariant holonomy
\begin{equation}
\mathcal U_\gamma
=
\exp\left(i\oint_\gamma \mathcal A\right).
\label{eq:u1-holonomy}
\end{equation}
In the electromagnetic Aharonov--Bohm realization this is conventionally written
\begin{equation}
\mathcal U_{\mathrm{AB}}
=
\exp\left(
\frac{i q}{\hbar}
\oint_\gamma A_\mu\,\dd x^\mu
\right),
\label{eq:ab-holonomy}
\end{equation}
up to the usual convention-dependent factors such as $c$ in nonrelativistic unit systems \cite{AharonovBohm1959}.  The relevant structural fact for the present programme is not a local force but a loop phase generated by a connection.

Berry phase has the same mathematical form: it is a $U(1)$ connection holonomy over parameter space \cite{Berry1984,Simon1983}.  Therefore the following statement is stronger than a loose analogy but weaker than a physical field identification.

\begin{proposition}[Gauge-geometric equivalence class]
\label{prop:u1-equivalence}
Aharonov--Bohm phase and Berry phase are both realizations of $U(1)$ connection holonomy.  If their dimensionless loop phases equal $\pi$ modulo $2\pi$, their holonomies both equal $-1$.
\end{proposition}

\begin{proof}
Both reduce to $\exp(i\theta)$ for a dimensionless loop phase $\theta$.  At $\theta=\pi\pmod{2\pi}$ the value is $-1$.
\end{proof}

The statement that the AB sector and the global radial sector below are two manifestations of one TIR ``potential tension'' is not established by this proposition and is recorded separately as a model postulate.

\section{Hubble normalization}
The Hubble parameter provides an inverse-time scale.  At the present epoch write $H_0$.  A corresponding length scale is
\begin{equation}
L_H=\frac{c}{H_0}.
\label{eq:hubble-length}
\end{equation}
The historical velocity--distance relation motivating the Hubble scale is reviewed from the original observational setting in \cite{Hubble1929}.  For the present construction we use only the dimensionless radial normalization
\begin{equation}
r_H(L)
=
\frac{H_0L}{c},
\label{eq:hubble-normalized-radius}
\end{equation}
for which
\begin{equation}
r_H(L_H)=1.
\end{equation}

\begin{warning}
Equation \eqref{eq:hubble-normalized-radius} is a dimensional normalization identity.  Interpreting $r_H=1$ as the literal boundary of the compactified zeta geometry, or interpreting $H_0$ as the same physical potential as the electromagnetic AB connection, is \emph{not} a theorem of cosmology or gauge theory.
\end{warning}

\section{Inverse--holonomy postulates for v0.2}
The new model layer is therefore stated explicitly rather than smuggled into the exact lemmas.

\begin{assumption}[Inverse--holonomy chart]
\label{ass:inverse-holonomy}
The physically relevant zero-state chart is centred at $s=1/2$, admits the reciprocal duality $u\mapsto1/u$, and is represented at finite radius by a disk compactification whose centre is $u=0$ and whose asymptotic region approaches the boundary.
\end{assumption}

\begin{assumption}[Potential-tension correspondence]
\label{ass:potential-tension}
After all quantities are made dimensionless, the local $U(1)$ loop phase represented by the Aharonov--Bohm/Berry holonomy sector and the global Hubble radial normalization are treated as two coordinates of one TIR potential-tension architecture.
\end{assumption}

\begin{assumption}[Holonomic zero closure]
\label{ass:holonomic-zero}
A canonical zero state closes with gauge-invariant loop phase
\begin{equation}
\theta_\gamma=\pi\pmod{2\pi},
\end{equation}
so that the corresponding $U(1)$ holonomy is $-1$.
\end{assumption}

Assumption \ref{ass:holonomic-zero} is compatible with the already proved Berry and two-channel cancellation results, but it does not yet derive the zero-state representation of Assumption \ref{ass:zero-state}.

\section{Claim promotions and remaining obstruction}
The v0.2 layer supports the following promotions:
\begin{enumerate}
\item \textbf{NEW THEOREM:} the centred zeta involution is $\widetilde J(u)=-\overline u$.
\item \textbf{NEW THEOREM:} reciprocal inversion commutes with $\widetilde J$.
\item \textbf{NEW THEOREM:} the critical axis is invariant as a set under $u\mapsto1/u$.
\item \textbf{NEW THEOREM:} the finite disk compactification places $u=0$ at the centre and $|u|=\infty$ at the boundary limit.
\item \textbf{PARTIAL PROMOTION OF C-006:} $p=\Ree s$ is the unique affine endpoint-preserving strip-to-binary-simplex coordinate.  Its use as the actual zero-state population remains unproved.
\item \textbf{STRUCTURAL PROMOTION:} Berry and Aharonov--Bohm phases are placed in the same mathematical $U(1)$ holonomy class.  Their physical potentials are not identified.
\end{enumerate}

The following claims do \emph{not} advance:
\begin{enumerate}
\item C-007, that every nontrivial zero is a genuine two-branch cancellation state, remains an \status{OPEN GAP};
\item pointwise self-duality of every zero remains unproved;
\item C-008, the Riemann hypothesis, remains \status{OPEN}.
\end{enumerate}

The mathematical centre of the next step is now narrower.  One must derive, from zeta data rather than impose, a faithful zero-state object for which the reciprocal symmetry, $U(1)$ holonomy, and affine strip coordinate jointly force local rather than merely pairwise balance.
